\documentclass[11pt]{amsart}

\usepackage[T1]{fontenc}
\usepackage{lmodern}
\usepackage{microtype}
\usepackage{amsmath,amssymb,amsthm}
\usepackage[hidelinks]{hyperref}

\hypersetup{
  pdftitle={Counterexamples to Lew's H-Free Laplacian Conjecture}
}

\newtheorem{theorem}{Theorem}
\newtheorem{lemma}[theorem]{Lemma}
\newtheorem{proposition}[theorem]{Proposition}
\newtheorem{corollary}[theorem]{Corollary}
\theoremstyle{remark}
\newtheorem{remark}[theorem]{Remark}

\newcommand{\eps}{\varepsilon}
\newcommand{\ex}{\operatorname{ex}}

\title[Counterexamples to Lew's conjecture]
{Counterexamples to Lew's
\texorpdfstring{$H$}{H}-Free Laplacian Conjecture}

\date{}

\subjclass[2020]{05C50, 05C35}
\keywords{Laplacian eigenvalue sums, forbidden subgraphs, Tur\'an numbers,
Brouwer-type inequalities}

\begin{document}

\begin{abstract}
For a graph $G$, let
\[
\eps_k(G)=\sum_{i=1}^k\lambda_i(L(G))-|E(G)|,
\]
where the Laplacian eigenvalues are ordered nonincreasingly. Lew conjectured
that every $H$-free graph $G$ with $|V(G)|\ge k$ satisfies
$\eps_k(G)\le \ex(k+1;H)$. We disprove the conjecture as stated: the path
$P_4$ is triangle-free, but
\[
\eps_2(P_4)=1+\sqrt2>2=\ex(3;K_3).
\]
This example is minimal in both $k$ and $|V(G)|$ when the forbidden graph has
an edge and no isolated vertices. More generally, every graph $H$ with
$\delta(H)\ge2$ admits a connected $H$-free counterexample, and for every
$r\ge3$ and $k\ge r-1$ there is a connected $K_r$-free counterexample at
parameter $k$.
\end{abstract}

\maketitle

\section{The smallest counterexample}

All graphs are finite and simple, and $H$-free means containing no copy of
$H$ as a subgraph. Let
\[
\lambda_1(L(G))\ge\cdots\ge\lambda_{|V(G)|}(L(G))=0
\]
be the Laplacian eigenvalues of $G$, and let $\ex(n;H)$ be the maximum number
of edges in an $n$-vertex $H$-free graph. Conjecture~7.1 of
Lew~\cite{Lew} asserts that
\begin{equation}\label{eq:lew}
  \eps_k(G)\le \ex(k+1;H)
\end{equation}
for every $k\ge1$, every graph $H$, and every $H$-free graph $G$ with
$|V(G)|\ge k$. Verbatim, Conjecture~7.1 of \cite{Lew} reads: ``Let $k\ge1$, and
let $H$ be a graph. Then, for every $H$-free graph $G=(V,E)$ with $|V|\ge k$,
$\eps_k(G)\le\ex(k+1;H)$.'' It is this statement, with no restriction on $H$
beyond being a graph, that we refute.

We write $P_j$ for the path on $j$ vertices.

\begin{theorem}\label{thm:p4}
The inequality~\eqref{eq:lew} is false. Specifically, for $k=2$, $H=K_3$,
and $G=P_4$,
\[
  \eps_2(G)=1+\sqrt2>2=\ex(3;K_3).
\]
\end{theorem}

\begin{proof}
A direct calculation gives
\[
\det(tI-L(P_4))=t(t-2)(t^2-4t+2),
\]
so the Laplacian eigenvalues of $P_4$ are
\[
2+\sqrt2,\qquad 2,\qquad 2-\sqrt2,\qquad 0.
\]
Since $P_4$ has three edges,
\[
\eps_2(P_4)=(2+\sqrt2)+2-3=1+\sqrt2.
\]
The graph $P_4$ is triangle-free, while Mantel's theorem gives
$\ex(3;K_3)=2$.
\end{proof}

\section{Pendant amplification}

For a graph $F$ and an integer $N\ge1$, let $F^{\langle N\rangle}$ be the
graph obtained by attaching $N$ new pendant vertices to each vertex of $F$.

\begin{lemma}\label{lem:pendant}
Let $F$ be a connected graph with $k\ge2$ vertices, $e$ edges, and Laplacian
eigenvalues $\mu_1,\ldots,\mu_k$. Put $G=F^{\langle N\rangle}$. Then the
Laplacian spectrum of $G$ consists of $1$ with multiplicity $k(N-1)$ and,
for each $i$, two numbers $\alpha_i,\beta_i$ that are the roots of
\begin{equation}\label{eq:quadratic}
  t^2-(N+1+\mu_i)t+\mu_i=0
\end{equation}
ordered so that $0\le\alpha_i<1<\beta_i$. Moreover,
\begin{equation}\label{eq:pendant-bound}
  \eps_k(G)
  =k+e-\sum_{i=1}^k\alpha_i
  >k+e-\frac{2e}{N+1}.
\end{equation}
\end{lemma}

\begin{proof}
For each vertex of $F$, the vectors supported on its $N$ pendant neighbors
and having coordinate sum zero give the eigenvalue $1$. Their total dimension
is $k(N-1)$.

Let $x$ be a Laplacian eigenvector of $F$ with eigenvalue $\mu$. On vectors
whose core coordinates are $ax$ and whose pendant coordinates adjacent to the
$i$th core vertex are all $bx_i$, the Laplacian acts on $(a,b)$ through
\[
\begin{pmatrix}
\mu+N&-N\\
-1&1
\end{pmatrix}.
\]
Its characteristic polynomial is \eqref{eq:quadratic}. The pendant-difference
vectors of the previous paragraph vanish on the core and have coordinate sum
zero on each pendant block, whereas the vectors just described are constant on
each pendant block; hence the two families are orthogonal, and in particular
linearly independent. Since $k(N-1)+2k=k(N+1)=|V(G)|$, together they form a
basis of $\mathbb{R}^{V(G)}$, so the numbers listed above are the whole
Laplacian spectrum of $G$.

For $\mu>0$, the polynomial in \eqref{eq:quadratic} is positive at $0$ and
negative at both $1$ and $N+1$. Hence its smaller root lies in $(0,1)$ and
its larger root exceeds $N+1$. For $\mu=0$, the roots are $0$ and $N+1$.
Thus the $k$ largest Laplacian eigenvalues of $G$ are the $\beta_i$.

Since $\alpha_i+\beta_i=N+1+\mu_i$ and $\sum_i\mu_i=2e$,
\[
\sum_{i=1}^k\beta_i
=k(N+1)+2e-\sum_{i=1}^k\alpha_i.
\]
The graph $G$ has $kN+e$ edges, which gives the identity in
\eqref{eq:pendant-bound}. Finally,
$\alpha_i\beta_i=\mu_i$ and $\beta_i>N+1$ whenever $\mu_i>0$, so
\[
\sum_{i=1}^k\alpha_i
<\frac{\sum_i\mu_i}{N+1}
=\frac{2e}{N+1}.
\]
\end{proof}

\begin{proposition}\label{prop:amplification}
Let $H$ satisfy $\delta(H)\ge2$, and let $F$ be a connected $H$-free graph
with $k\ge2$ vertices and $e$ edges. Set
\[
  d=k+e-\ex(k+1;H).
\]
If $d>0$ and $2e/(N+1)\le d$, then $F^{\langle N\rangle}$ is a connected
$H$-free counterexample to \eqref{eq:lew} at parameter $k$.
\end{proposition}

\begin{proof}
A copy of $H$ cannot use a pendant vertex, because every vertex of $H$ has
degree at least two. Hence $F^{\langle N\rangle}$ is $H$-free, and it is
connected because $F$ is connected. Lemma~\ref{lem:pendant} gives
\[
\eps_k\bigl(F^{\langle N\rangle}\bigr)
>k+e-\frac{2e}{N+1}
\ge k+e-d
=\ex(k+1;H).
\]
\end{proof}

\begin{corollary}\label{cor:all-H}
For every graph $H$ with $\delta(H)\ge2$, there is a connected $H$-free
counterexample to \eqref{eq:lew}.
\end{corollary}

\begin{proof}
Let $h=|V(H)|$, set $k=h-1$, and take $F=K_k$. Since $\delta(H)\ge2$ while
degrees in a simple graph are at most $h-1$, we have $h\ge3$; hence $k=h-1\ge2$
and
$e=\binom{k}{2}\ge1$, as required by Lemma~\ref{lem:pendant} and
Proposition~\ref{prop:amplification} and so that $N=2e\ge1$. Then $F$ is
connected and $H$-free. Moreover,
\[
k+e=\binom{h}{2}
\quad\text{and}\quad
\ex(h;H)\le\binom{h}{2}-1,
\]
because $K_h$ contains $H$. Thus $d\ge1$. Since $2e/(2e+1)<1$, taking
$N=2e$ in Proposition~\ref{prop:amplification} proves the claim.
\end{proof}

\begin{corollary}\label{cor:complete}
For every $r\ge3$ and every $k\ge r-1$, there is a connected $K_r$-free
graph $G$ such that
\[
  \eps_k(G)>\ex(k+1;K_r).
\]
\end{corollary}

\begin{proof}
Let $F=T_{r-1}(k)$ be the balanced complete $(r-1)$-partite graph on $k$
vertices, and write
\[
q=|E(F)|=\ex(k;K_r).
\]
By Tur\'an's theorem, adding one vertex to a smallest part gives
\[
\ex(k+1;K_r)-\ex(k;K_r)
=k-\left\lfloor\frac{k}{r-1}\right\rfloor.
\]
Consequently,
\[
k+q-\ex(k+1;K_r)
=\left\lfloor\frac{k}{r-1}\right\rfloor\ge1.
\]
The graph $F$ is connected because $k\ge r-1$. Since
$2q/(2q+1)<1$, Proposition~\ref{prop:amplification}, applied with
$H=K_r$ and $N=2q$, proves the claim.
\end{proof}

\begin{remark}\label{rem:size}
The violation produced above is small. By the identity in
\eqref{eq:pendant-bound}, with $\alpha_i$ as in Lemma~\ref{lem:pendant} and
$F=T_{r-1}(k)$, the graph $G=T_{r-1}(k)^{\langle2q\rangle}$ satisfies
\[
\eps_k(G)-\ex(k+1;K_r)
=\left\lfloor\frac{k}{r-1}\right\rfloor-\sum_{i=1}^k\alpha_i
<\left\lfloor\frac{k}{r-1}\right\rfloor,
\]
so the excess is $O(k)$, while $\ex(k+1;K_r)=\Theta(k^2)$. The reason is
visible in \eqref{eq:pendant-bound}: pendant amplification gains less than $k$
over $|E(F)|$, whereas passing from $k$ to $k+1$ vertices raises the Tur\'an
number by $k-\lfloor k/(r-1)\rfloor$, falling short of $k$ by exactly
$\lfloor k/(r-1)\rfloor$. Thus it is the exact inequality \eqref{eq:lew} that
fails, and not its asymptotic form: the counterexamples are consistent with the
bound $\eps_k(G)\le\ex(k;H)+(4k-2)\sqrt k$, valid for every $H$-free graph $G$
and every $1\le k\le|V(G)|$, recorded in \cite[Section~7]{Lew}, and we
do not know whether \eqref{eq:lew} holds up to a multiplicative $1+o(1)$ or an
additive $o(k^2)$ error.
\end{remark}

\section{Minimality}

\begin{proposition}\label{prop:minimality}
Suppose that $H$ has at least one edge and no isolated vertices. Then
\eqref{eq:lew} holds when $k=1$, and it also holds whenever
$|V(G)|\le k+1$. Hence every counterexample under these assumptions satisfies
\[
k\ge2,
\qquad
|V(G)|\ge k+2\ge4.
\]
The example in Theorem~\ref{thm:p4} attains both lower bounds.
\end{proposition}

\begin{proof}
First let $k=1$. If $G$ is edgeless, then $\eps_1(G)=0$. Otherwise choose a
connected component $C$ with
$\lambda_1(L(C))=\lambda_1(L(G))$, and put $t=|V(C)|$. Since $C$ is
connected,
\[
t\le |E(C)|+1\le |E(G)|+1.
\]
Also $\lambda_1(L(C))\le t$: on the orthogonal complement of the all-ones
vector,
\[
L(C)+L(\overline C)=tI.
\]
Therefore $\eps_1(G)\le1$. If $H=K_2$, every $H$-free graph is edgeless. If
$H\ne K_2$, then $\ex(2;H)=1$. Thus \eqref{eq:lew} holds for $k=1$.

Now let $n=|V(G)|\le k+1$. Since $n\ge k$, either $k=n$ or $k=n-1$.
In both cases the sum of the $k$ largest Laplacian eigenvalues is
$2|E(G)|$, so
\[
\eps_k(G)=|E(G)|.
\]
If $n=k+1$, the definition of $\ex(k+1;H)$ gives the desired bound. If
$n=k$, add one isolated vertex to $G$. Because $H$ has no isolated vertices,
the resulting $(k+1)$-vertex graph remains $H$-free, and the same bound
follows.
\end{proof}

\begin{remark}
The no-isolated-vertices hypothesis in Proposition~\ref{prop:minimality} is
necessary. For
\[
H=K_2\cup K_1,\qquad G=K_2,\qquad k=2,
\]
one has $\eps_2(G)=1>0=\ex(3;H)$. In particular, since \eqref{eq:lew} is
asserted for every graph $H$, the conjecture in its literal form already fails
for $G=K_2$ at $k=2$, with only two vertices. The minimality assertion of
Proposition~\ref{prop:minimality} concerns forbidden graphs with an edge and
no isolated vertices, and it is in that regime that $P_4$ is smallest.
\end{remark}

\begin{thebibliography}{9}

\bibitem{Lew}
A.~Lew,
\emph{An approximate version of Brouwer's Laplacian conjecture},
arXiv:2601.17575v1 [math.CO], 2026.

\end{thebibliography}

\end{document}